\section{Related Work}
\label{sec:related}

\subsection{Spatial and spatiotemporal reasoning in MLLMs}
SpatialVLM studies how to equip image-language models with quantitative spatial knowledge, while SpatialRGPT grounds spatial answers in regions and SpatialBot extends the input toward RGB-D geometry~\citep{chenSpatialVLMEndowingVisionLanguage2024,chengSpatialRGPTGroundedSpatial2024,caiSpatialBotPreciseSpatial2025}. Scene-LLM projects camera- and world-coordinate 3D scene features into an LLM using staged alignment and instruction tuning~\citep{fuSceneLLMExtending2025}. These approaches make depth, position, and object relationships more accessible to language generation. Video adds another difficulty: an object can change position while the camera moves, and its apparent displacement may not equal its motion in the world. VSI-Bench, OmniSpatial, and ViSTR-Bench document persistent shortcomings on spatial and spatiotemporal queries~\citep{yangThinkingSpaceHow2025,jiaOmniSpatialComprehensiveSpatial2026,liViSTRBenchCanMLLMs2026}. Dynamic-scene methods attempt to narrow this gap. VG-LLM uses geometry-oriented video supervision; VLM-3R links reconstruction and language learning; MLLM-4D uses instruction tuning for time-varying spatial relationships~\citep{zhengLearningVideos3D2025,fanVLM3RVisionLanguageModels2026,yinMLLM4DVisualbasedSpatialTemporal2026}. Their representation, supervision, and test questions differ, so their results do not directly answer whether a language model can perform a numerical motion query when world-frame tracks are already available. Our setting isolates that interface: it gives the MLLM reconstructed object and camera trajectories and asks it to read metric positions and elapsed times from them. This is complementary to methods that infer geometry from video. It also makes the dependence on the reconstruction explicit; strong performance here cannot by itself establish accurate motion recovery from raw pixels. The frames-only arm in our experiments tests the practical value of this extra input on the same questions, while donor-token and question-only controls ask whether the answer actually depends on the supplied trajectory values.

\subsection{Dynamic 3D reconstruction and object tracks}
A broad line of work estimates correspondences and geometry across a moving camera and moving objects. Stereo4D mines world-consistent, pseudo-metric point tracks from internet stereo video and provides the reconstructed trajectories on which our data and reference answers are based~\citep{jinStereo4DLearningHow2025}. Methods such as TAPIP3D and SpatialTrackerV2 pursue persistent 3D point tracking under their respective input assumptions~\citep{zhangTAPIP3DTrackingAny2025,xiaoSpatialTrackerV23DPoint2025}. Track4World predicts world-centric dense 3D tracks from monocular video~\citep{luTrack4WorldFeedforwardWorldcentric2026}; our broader pipeline can use it when released tracks are absent, but every downstream TrajQA-B result in this paper uses Stereo4D's provided tracks. This distinction matters because a better tracker could change both the input and the reconstruction-derived target. Dense tracks are also not an object-level language interface: they can contain many redundant surface points, change visibility across time, and mix rigid camera motion with independently moving objects. We therefore use masks to attribute tracks to instances, select representative points for their geometry or semantic parts, and retain a separate camera trajectory. The object inventory and part assignment may themselves be imperfect, which is why the benchmark gates retain only cases whose answers agree when recomputed from the serialized input. Our contribution begins after motion reconstruction: it organizes already estimated tracks into a compact, queryable sequence. We do not claim that this organization improves the underlying 3D estimates, and our evaluation inherits the reconstruction's scale, occlusion, and identity errors. This places our method between 4D reconstruction and language reasoning rather than presenting it as a replacement for either.

\subsection{Trajectory representations and language interfaces}
Continuous trajectories pose a representation choice: uniform frame sampling preserves a simple temporal grid but spends tokens on nearly redundant positions and may skip a brief change; quantization and learned codes can reduce sequence length but may obscure metric values. One Trajectory, One Token organizes video features around panoptic sub-object tracks and connects the encoder to VideoLLM QA~\citep{zhengOneTrajectoryOneToken2025}. TrajTok adapts spatial tokenization of GPS traces~\citep{xiongTrajTokAdaptiveSpatial2026}, while LLaMo represents human motion with continuous autoregressive latents~\citep{liLLaMoScaling2026}. MotionLM and AutoTraces use language-model mechanisms for motion forecasting~\citep{seffMotionLMMultiAgentMotion2023,wangAutoTracesAutoregressiveTrajectory}. These works mostly encode or generate motion, whereas ours supplies reconstructed metric motion to answer queries from recoverable numeric records. We fit each visible keypoint run with a bounded-error piecewise curve, allocate more nodes where the path changes, and serialize on-curve samples with object, part, physical-time, and event fields. The representation also carries the normalization parameters, so a model can in principle reconstruct metre-scale positions rather than relying on relative motion codes. A learned projector maps each record to the language model's embedding space without changing its visual encoder; relations among records are learned inside the language model. This direct interface is related to trajectory tokenization in its budget goal, but differs in its emphasis on object identity, visibility gaps, camera motion, and verifiable metric answers. Our current intervention experiments establish dependence on correct token content as a whole. They do not yet isolate the necessity of adaptive fitting, part labels, or event flags, which remain specific component-level hypotheses for further ablation.
